\section{Yahweh’s Template}

\begin{quotationx}
For as lightning that comes from the east is visible even in the west, so will be the coming of the Son of Man.

In the madness of materialism the West delivers its great thinkers to the graveyard of thoughts, and tramples those in the dirt, who wish to abjure this madness with strong and holy words.
\flright{\textsc{Prince V. F. Odoyevsky}}
\end{quotationx}


A common understanding today is that Christianity is inherently Semitic, while the Gnostics attempted to embrace a deeper spirituality more in tune with the dramatic nature of man's fall. There is some truth to this, both emotionally \& metaphysically; we have posted Dugin's essay\footnote{\url{http://www.arctogaia.com/public/eng/gnostic.htm}} on Gnosticism, which should resonate with any denizen of the brave new world. Still, the Church (which is surrounded by a ``halo of hatred'', in Chesterton's words) is more complex than this over-simplification allows. The Gnostic-Christian debate has all the hallmarks of schizophrenia, a split in the mind of Mediterranean men, playing out in such sad tragedies as the Albigensian Crusades, in which Christendom made war upon her own bosom in an attempt to purge the self of an evil, perceived or not. We cannot de-Hellenize the Church or the Faith, yet the Faith sleeps uneasily with its Hellenistic past.

History at this level presents a Gordian knot; the same levels of thinking which produced such an impasse, will not resolve the impasse. Hence, we turn to esoteric study.

Reputedly, the diaries of Clement\footnote{\url{http://jacksonsnyder.com/arc/bible-news/apostolia/pdf/ALL-CLEMENT-6x9-120707-embedded.pdf}} -- or here\footnote{\url{http://www.scribd.com/doc/13561289/-The-Authentic-Peter-The-Preaching-of-Simeon-Kefa-from-the-Journal-of-T-Flavius-Clemens-Clement-The-Clementine-Recognitions-Recognitions-of-Cl}} -- have been translated out from their oblivion in the \emph{Ante-Nicene Fathers}. In this work, there is an argument between Bar Cepha (Cephas, named Peter) and Shimon concerning the nature of the Creator. Shimon is a Gnostic who believes that he has a deeper truth than the Torah, coming from beyond time and space, tainted by the Architect.

During the debate (in which Cephas ``the Rock'' interestingly makes the point that reason \& faith are necessary for firm he does not say true knowledge of God, rather than mere ``belief'' alone) Saint Peter admits that he \emph{does} in fact know \emph{where and whence} eternal spirits came, but that to discuss their origin would do Shimon no good, as Shimon has already begun with the false premise that Yahweh is ignorant of the true Father of Lights to whom the spirits belong. To believe that Yahweh ignorantly or in delusion created the fallen material bodies \& physical world is to be so utterly wrong that it excludes one from the mysteries.

This is entirely consistent with the apostolic tradition. It is little known or appreciated how secretive the early church was\footnote{\url{http://www.earlychristianwritings.com/text/hippolytus-christ.html}}:

\begin{quotex}
By these, too, you will be able to silence those who oppose and gainsay the word of salvation. Only see that you do not give these things over to unbelieving and blasphemous tongues, for that is no common danger. But impart them to pious and faithful men, who desire to live holily and righteously with fear. For it is not to no purpose that the blessed apostle exhorts Timothy, and says, ``O Timothy, keep that which is committed to thy trust, avoiding profane and vain babblings, and oppositions of science falsely so called; which some professing have erred concerning the faith.'' And again, ``Thou therefore, my son, be strong in the grace that is in Christ Jesus. And the things that thou hast heard of me in many exhortations, the same commit thou to faithful men, who shall be able to teach others also.'' If, then, the blessed (apostle) delivered these things with a pious caution, which could be easily known by all, as he perceived in the spirit that ``all men have not faith,'' how much greater will be our danger, if, rashly and without thought, we commit the revelations of God to profane and unworthy men?
\end{quotex}


So Cephas is following Christ's method, which would be well known -- teach in parables, do not explain to Pharisees.

Still it is unclear if Cephas' objection is entirely merely that the wrong is done to ``the Creator'' -- after all, as we will see, he is willing to wrangle a Pascalian wager. Perhaps, the more immediate problem is viewing \emph{matter} as \emph{inherently} evil. If one has to ascend, then it is more important to be right about \emph{matter} than about God, since matter is what you are closest to. A matter which is purely evil would preclude (as we shall see) any higher knowledge as well. Hence, in esoteric studies, the need to begin with parables \& with baby steps (as Christ does). To transcend, start to pay attention.

For metaphysically, Cephas' more basic argument is that Gnosticism does not ``resolve'' anything about the problem of evil except by sleight of hand -- if the True Father of Lights can aid us, he knows of us, and this means he knows of our plight. This would entail his foreknowledge of what Yahweh did, which entails the same dilemma which the Gnostics pose regarding Yahweh, \& for which they indict him of being ``lesser''. Cephas, in short, faced with the prospect of doubting all the known senses, would rather leave the conundrum in Yahweh's court -- the Father of Lights is more likely to have mercy (in case Yahweh is truly ``lesser''). In other words, \emph{we have here the basics of the argument of infinite regression \textbf{combined} with Pascal's wager} -- Cephas is betting that the Torah does have a higher meaning, but that this meaning is not completely unconsonant with the literal, childlike meaning which is necessary to begin understanding the structure of the Cosmos. Most Christian orthodoxies unconsciously acknowledge, point to, and tend towards this truth at their subtle core; or more precisely, there is an element within them that does, at the least. Provided the Church does not fall into materialistic moralizing of the most grotesque sort, it should always has this drift. We know this debate as Law vs. Grace, but it takes many forms in Christianity, because Gnosticism was once part of her esoteric core, before it was ejected (perhaps self-ejected?).

As an example, this point is similar to the Pope's recent claim in his speech on Pseudo-Dionysius\footnote{\url{https://gornahoor.net/?p=2332}} -- there is a profound innate similarity between the language of the faithful (Calvin's method of discussing God) \& the knowledge that goes beyond all words.

If we accept the Cephas-Mark-Clement connection in Alexandria, we can place esoteric Christianity in close proximity with almost all of the currents present in the Classical world (some claim that the Egyptian connection is a red herring, \& the Eleusianian mysteries\footnote{\url{http://www.scribd.com/doc/28988441/Katabasis-The-Eleusinian-Mysteries-And-the-Early-Christian-Church}} in fact are the real link\footnote{\url{http://books.google.com/books?id=zXPvqI9fz7QC&pg=PA124&lpg=PA124&dq=eleusinian+mysteries+christian+liturgy&source=bl&ots=MqdNdP8N35&sig=YlUhdPpZXM8GSTfze4mDBU3mnKE&hl=en&ei=Q2cSTrroK-Tr0gH2ncmODg&sa=X&oi=book_result&ct=result&resnum=1&ved=0CBcQ6AEwADgK\#v=onepage&q=eleusinian\%20mysteries\%20christian\%20liturgy&f=false}}).

Here is Cicero:

\begin{quotex}
For among the many excellent and indeed divine institutions which your Athens has brought forth and contributed to human life, none, in my opinion, is better than those mysteries. For by their means we have been brought out of our barbarous and savage mode of life and educated and refined to a state of civilization; and as the rites are called ``initiations,'' so in very truth we have learned from them the beginnings of life, and have gained the power not only to live happily, but also to die with a better hope.
\end{quotex}


It is well known that the Roman nobility contributed many of their class to the ranks of bishops (eg, in Aquitainia), and so we can concretely observe Christianity incorporate the civilizing mission of Rome.

To reiterate, Christianity is inseparable from -but unreconciled to -- her Hellenic past. She cannot either totally forget, or completely remember.

In trying to reconstruct historical trajectories without ``the mysteries'', we encounter great difficulty. The legend of Andrew going to Russia\footnote{\url{http://en.wikipedia.org/wiki/Saint\_Andrew}} (Cyril and Methodius are said to have found the Gospels already translated there) makes sense, \& might put us closer to the point of reconstructing the seminal explosion of the Faith in the early centuries.

It would be easier if we had ``the template''. Fortunately, the Scriptures themselves (as the quotations by Hippolytus demonstrate) are for more rewarding in this way than most suppose.

Jesus began with the Beatitudes (Matthew 5:3-12)

\begin{itemize}


\item the poor in spirit: for theirs is the kingdom of heaven\footnote{\url{http://en.wikipedia.org/wiki/Kingdom\_of\_God}}. (5:3\footnote{\url{http://en.wikipedia.org/wiki/Matthew\_5:3}})

\item they that mourn: for they shall be comforted. (5:4\footnote{\url{http://en.wikipedia.org/wiki/Matthew\_5:4}})

\item the meek: for they shall inherit the earth. (5:5\footnote{\url{http://en.wikipedia.org/wiki/Matthew\_5:5}})

\item they which do hunger and thirst after righteousness: for they shall be filled. (5:6\footnote{\url{http://en.wikipedia.org/wiki/Matthew\_5:6}})

\item the merciful\footnote{\url{http://en.wikipedia.org/wiki/Works\_of\_Mercy}}: for they shall obtain mercy\footnote{\url{http://en.wikipedia.org/wiki/Mercy}}. (5:7\footnote{\url{http://en.wikipedia.org/wiki/Matthew\_5:7}})

\item the pure in heart: for they shall see God. (5:8\footnote{\url{http://en.wikipedia.org/wiki/Matthew\_5:8}})

\item the peacemakers: for they shall be called the children of God. (5:9\footnote{\url{http://en.wikipedia.org/wiki/Matthew\_5:9}})

\item they which are persecuted for righteousness' sake: for theirs is the kingdom of heaven. (5:10\footnote{\url{http://en.wikipedia.org/wiki/Matthew\_5:10}})
\end{itemize}


If one compares the Beatitude progression of poverty in spirit, merciful, etc. (8 steps) to passages in Scripture from 2 Peter (faith, self-control, knowledge, etc. leading to love in 8 steps), \& match them against St. Paul's progression of faith, hope, then love (as outlined in Mouravieff's \emph{Gnosis}), one gets a very close match.

\begin{quotex}
For this very reason, make every effort to add to your faith goodness; and to goodness, knowledge; \textsuperscript{6} and to knowledge, self-control; and to self-control, perseverance; and to perseverance, godliness; \textsuperscript{7} and to godliness, mutual affection; and to mutual affection, love. \textsuperscript{8} For if you possess these qualities in increasing measure, they will keep you from being ineffective and unproductive in your knowledge of our Lord Jesus Christ.
\textsuperscript{9} But whoever does not have them is nearsighted and blind, forgetting that they have been cleansed from their past sins.
\end{quotex}


Also confirming this by the \emph{Shepherd of Hermas} --

\begin{quotex}
Hermas 22 1 When then I ceased asking her concerning all these things, desirous of beholding, I was greatly rejoiced that I should see it. 2 She looked upon me, and smiled, and she saith to me, Seest thou seven women round the tower? 3 I see them, lady, say I. 4 This tower is supported by them by commandment of the Lord. 5 Hear now their employments. 6 The first of them, the woman with the strong hands, is called Faith; 7 through her are saved the elect of God. 8 And the second, that is girded about and looketh like a man, is called continence; 9 she is the daughter of Faith. 10 Whosoever then shall follow her, becometh happy in his life, for he shall refrain from all evil deeds, believing that, if he refrain from every evil desire, he shall inherit eternal life. 11 And the others, lady, who be they? 12 They are daughters one of the other. 13 The name of the one is Simplicity, of the next, knowledge, of the next, Guilelessness, of the next, Reverence, of the next, Love. 14 When then thou shalt do all the works of their mother, thou canst live. 15 I would fain know, lady, I say, what power each of them possesseth. 16 Listen then, saith she, to the powers which they have. 17 Their powers are mastered each by the other, and they follow each other, in the order in which they were born. 18 From Faith is born Continence, from Continence Simplicity, from Simplicity Guilelessness, from Guilelessness Reverence, from Reverence knowledge, from knowledge Love. 19 Their works then are pure and reverent and divine. 20 Whosoever therefore shall serve these women, and shall have strength to master their works, shall have his dwelling in the tower with the saints of God.
\end{quotex}


-- we can begin to discern a Christian ``Tao'' which was incorporating Greek mysteries and Roman rites in an effort to embody the ``Kingdom of God''. \emph{It was this ideal, and not that of cheap grace or individual salvation, which fired the hearts \& minds of Europeans during the Crusades and long after}. It was this which went to the four corners -- so if we are looking for this template, we may seek \& find. The Faith is self-awareness \& fullness of this embodied in a concrete civilization. The historical trajectories are suggestive clues, but we have to follow them, and expect to find them.

Again, the Faith represents in a way the balancing of the chakras of the universal Man, in the heart. The Faith (or can be) the apex of spiritual evolution for man. This places her at odds with those like Crowley who would seek to bypass the union of heaven and earth, or unbalance it. It is this struggle which defines the very headship of Christ; of His body and His many possessions we do not speak, but only of those called to be the lightning of God.

The power of the template to identify and predict the progress of the kingdom of God is astounding. This progression of knowledge applies at a civilizational level as well as a personal one. The two nation-civilizations most consciously working out the implications of revelation in the 1900s (for good or bad) were Germany \& Russia. The last century centers (basically in a negative way) on the German-Russian axis. \emph{The Crisis Of Civilization} (\emph{Anthroposophie auf der Kreuzung der okkult-politischen Bewegungen der Gegenwart}) by the excommunicated Russian Anthroposophist Gennady Bondarev is worth examining, along with Bondarev's \emph{Die Wartende Kultur}\footnote{\url{http://www.amazon.de/Die-wartende-Kultur-Esoterische-russischen/dp/3906712028}}, Prokofieff's \emph{Spiritual Origins of Eastern Europe}, \& Irina Gordienko's \emph{S.O. Prokofieff, Myth and Reality}. All of these works debate questions which Rudolf Steiner raised about dimensional human evolution, and where and how it would occur in relation to the world-cultures. The ideal German culture is supposed to come into expression via Idealism, Goetheanism, Anthroposophy, and Social Threefolding; the Slavic destiny (it is debated) instead amounts to bearing the Beast-State karma of world degeneration, thereby enabling other countries to hopefully gain ground\footnote{\url{http://www.altanthroinfo.9f.com/SpiritualEurope.htm}}. This is common ground in Steiner studies, apparently, and in related anthroposophy.\footnote{\url{http://en.wikipedia.org/wiki/New\_Galilee\_\%28Sixth\_Epoch\%29}}

For better or worse, the Church sought a harmony between exoteric tradition and esoteric mystery. Our task as traditionalists ought to be to enable the Church to better fulfill that aspiration (by understanding the process as she should do), without worrying too much about whether she wishes to acquiesce in all particulars. Part of the terrible challenge devastating the Church today is her refusal to ``scientifically'' or factually embrace a concrete and concretely demonstrated practical metaphysic to explain her esoteric doctrine she has made exoteric (see George Heart's\footnote{\url{http://www.amazon.com/Christianity-Dogmatic-Gnostic-Vivifying-Knowledge/dp/1412062284}} \emph{Dogmatic Faith \& Gnostic Vivifying Knowledge}). She has forgotten her own template, or dare not realize it. How could the Church not see that Hitler \& Stalin's advent presaged knowledge about the state of the Kingdom of God? Instead, she chose the ``darkness'' of refusing to see which Solzhenitsyn claimed had overtaken us at Nuremberg.

The last century was not the century without God -- it was the century which realized the inversion of the divine template in history, and its increasing amnesic treatment in a Church pursuing social gospel \& fundamentalism. Once the power of this template for identifying \& regenerating Christianity is realized, it will be possible to begin to discuss in all humility the real challenges of the Church, including ancient divisions with present consequences.

Yahweh, for instance, may be a particular manifestation or theophany of God which is transcended (Jesus speaks of His ``Father'') a particular version of this is precisely\footnote{\url{http://www.romanity.org/htm/rom.18.en.augustine\_unknowingly\_rejects\_the\_doctrine.01.htm}} the Eastern Orthodox dogma, in fact, where the transcendence is in a manner of speaking added to, or subsumed at the Transfiguration. There are, of course, other possibilities, including the idea that Yahweh is an angelic ruler who is deluded, but that matter is not \emph{inherently} evil, but is redeemed by its connection with spirit, however minute. However, this is not exclusive of the first possibility -- God may be ``drawing'' for children here.

\begin{quotex}
Formerly\footnote{\url{http://www.altanthroinfo.9f.com/SpiritualEurope.htm}} the Archangel Michael was the countenance of Jahve. Now he is the countenance of God... 
\end{quotex}

This is probably central to demoralizing or de-sentimentalizing the Church. In any case, the Gnostic emphasis was correct (flesh has to be transcended) while their understanding of Torah is wrong. In the Cephas debate, one senses that Cephas is dealing with someone who is getting steps out of order. If the Church had lost (or was unable to maintain or to teach) the proper steps as outlined in the Beatitude (combined with esoteric initiations), then the Gnostic misunderstanding and tragedy makes perfect sense. Further research into early Christianity \& the Greek or Russian connections may reveal more interesting facts and interpenetrations than we have endeavored to touch on here for now.

For more on Clement\footnote{\url{http://jacksonsnyder.com/arc/bible-news/apostolia/index.htm}}... 

\flrightit{Posted on 2011-07-09 by Logres}

\begin{center}* * *\end{center}

\begin{footnotesize}
\begin{sffamily}
\texttt{Theopathic on 2011-07-09 at 12:14 said:}

Forgive me, this was a truly splendid and theo-philosophically interesting and instructive communication, but I am writing an essay about Evola and cannot find, for the life of me, the actual source-text for Evola's citation of the ``Nitisara'' (4.4) on the semi-deific nature of ideal rulers in ``Revolt Against the Modern World.'' Can any one help out here with this haphazard and inelegant question?

\hfill

\texttt{logres on 2011-07-09 at 13:17 said:}

Granted, Faith = Poverty of Spirit, Goodness = Mourning, Knowledge = Meekness, Self-Control = Hungering/Thirsting for Righteousness is a kind of shift of tone, almost a progression within the progression, so that the first (Faith) is a product of the second (Poverty of Spirit). However, all the progressions move from Faith through Knowledge to Vision of God.

\hfill

\texttt{Heraklion on 2011-07-10 at 02:50 said:}

Inspiring words, Logres. Achieving the task of non-linearly, non-syncretically combining the Truth of Torah and Christ with the remnants of so-called ``pagans'' whose hearts still have God's Law imprinted on them (see Paul, Romans, of course) is a mind-numbingly difficult task... Certain better-natured Masons and Rosicrucians attempted such-like things in the past, but failed... A brotherhood between Godly Pagans and Godly Christians as a counterrevolutionary militia opposing the Kali Yuga is seriously exciting to at least imagine... \hfill

\texttt{logres on 2011-07-11 at 00:40 said:}

As Cologero has pointed out, a huge danger is betrayal from inside -- this is how occult power always work, and no doubt why such projects came to nothing in the past. Yet the biggest challenge (perhaps) is to manage to coordinate without achieving the very occult goal itself -- the attainment of a synthesis which combines the worst elements of both. Such might stand as a definition of ``modern times''.

\hfill

\texttt{logres on 2011-07-11 at 00:45 said:}

Theopathic,

It's in an easy on Junger, isn't it:

``In the Nitisara we are asked to explain how he who cannot dominate himself (his own manas) can dominate other men, and in the Arthaçâstra the exercise of royal functions is conceived as tapas, i.e. ascetism, ascetism of power. We might easily multiply references of this kind.''

If you are asking where Evola references it...

\url{http://euro-synergies.hautetfort.com/archive/2009/02/19/julius-evola-on-ernst-junger.html}


\hfill

\texttt{Count Cagliostro on 2011-07-11 at 02:54 said:}

``To reiterate, Christianity is inseparable from -but unreconciled to -- her Hellenic past. She cannot either totally forget, or completely remember''

Coming from Greece, I cannot but admire the wise choice of your words!


\hfill

\end{sffamily}
\end{footnotesize}

